 \subsection*{Introduction }

This is a the splash page of the documentation of the \char`\"{}core\char`\"{} S\+M\+S++, a set of C++ classes intended to provide a system for modeling complex, block-\/structured mathematical models (in particular, but not exclusively, single-\/real-\/objective optimization problems), and solving them via sophisticated, structure-\/exploiting algorithms (in praticular, but not exclusively, decomposition approaches and structured Interior-\/\+Point methods).

\subsection*{A General Description of S\+M\+S++ }

Current systems for modeling mathematical models typically falls into two opposite approaches\+:


\begin{DoxyItemize}
\item {\itshape General-\/purpose Algebraic Modeling Languages}, commercial ones like \href{https://ampl.com}{\tt A\+M\+PL} and \href{https://www.gams.co/}{\tt G\+A\+MS}, or open-\/source ones like \href{http://www.coliop.org}{\tt Coliop} and \href{http://zimpl.zib.de}{\tt Z\+I\+M\+PL}. These try to abstract away as much as possible the mathematical model from the underlying solver, thereby taking away from the user all the effort from interfacing with specific solvers tackling their idiosyncrasies. However, in order to do to they typically \char`\"{}flatten\char`\"{} all the structure available in the model, making it almost unreachable from the underlying solver. Some extensions have been proposed for specific forms of structure, mainly the block one, e.\+g. in the \href{https://www.maths.ed.ac.uk/ERGO/sml}{\tt Structured Modeling Language (S\+ML)}. An alternativa approach has been to re-\/construct the structure backwards from the \char`\"{}flat\char`\"{} model, as \href{http://www.or.rwth-aachen.de/gcg}{\tt Generic Column Generation (G\+CG)} does.
\item The alternative has typically been to {\itshape directly interface to a specific solver via its A\+P\+Is}; all modern general-\/purpose (and specialized) solvers, both commercial ones like \href{https://www.ibm.com/analytics/cplex-optimizer}{\tt Cplex}, \href{http://www.gurobi.com}{\tt Gu\+Ro\+Bi} and \href{https://www.mosek.com}{\tt M\+O\+S\+EK}, and open-\/source ones like \href{https://projects.coin-or.org/Cbc}{\tt Cbc} and \href{http://scip.zib.de}{\tt S\+C\+IP}, offer one or more A\+P\+Is, usually for different languages (among which, surely C/\+C++). This allows to directly address all the structure-\/exploiting features of the solver, if any, but directly explosed the user to its complexity and idiosyncrasies.
\item A somewhat \char`\"{}in the middle\char`\"{} approach is that of {\itshape modeling systems embedded in general purpose languages}, such as \href{https://projects.coin-or.org/FlopC++}{\tt Flop\+C++} and \href{https://github.com/coin-or/Rehearse}{\tt C\+O\+I\+N-\/\+Reharse} for C++, \href{https://pythonhosted.org/PuLP}{\tt Pu\+LP} and \href{http://www.pyomo.org}{\tt Pyomo} for Python, \href{https://github.com/JuliaOpt/JuMP.jl}{\tt Ju\+MP} for Julia and \href{https://yalmip.github.io}{\tt Y\+A\+L\+M\+IP} for Matlab. This allows to better deal with parts of the model that have specific structure and allow for specialized solution method, and thereby implement sophisticated structure-\/exploiting algorithms. However, at their hearth these systems typically still construct \char`\"{}flat\char`\"{} models, and have little explicit support for managing the structure in the model.
\end{DoxyItemize}

S\+M\+S++ tries a different approach, somehow extending the latter one but even more deeply rooted in programming languages (after all, mathematical optimization used to be known as \char`\"{}mathematical programming\char`\"{} for a good reason). It is a modeling system based on a set of C++ classes, where a model is represented as a \char`\"{}block\char`\"{}, containing some collections of \char`\"{}variables\char`\"{} and \char`\"{}constraints\char`\"{}, plus one \char`\"{}objective function\char`\"{}, possibly organized (recursively) into sub-\/blocks. The current distribution provides the corresponding {\itshape abstract} Block, Variable, Constraint and Objective classes, that try to establish the minimal possible interface of such a modeling system. In particular, Variable, Constraint and Objective make as little assumptions as possible on the actual form of the corresponding mathematical objects. However, a Real\+Objective class is already derived from Objective in order to al least define the crucial concept of an objective function with real values.

A first level of specification is then provided by the derived classes Col\+Variable (a single real Variable, possibly with integrality restrictions), and Row\+Constraint (a Constraint of the form \char`\"{}\+L\+H\+S $<$= ( single
real expression ) $<$= R\+H\+S\char`\"{}). Also, a general Function class is provided for expressing a very general real-\/valued function concept, together with its derived C05\+Function (providing linearizations, i.\+e., first-\/order information, but not necessarily continuous one) and C15\+Function (providing quadratic models, i.\+e., second-\/order information, but not necessarily continuous one). With this, F\+Row\+Constraint and F\+Real\+Objective are defined as Row\+Constraint and Real\+Objective taking a generic Function to implement them. Some basic aspects of the interface of Constraint, Objective and Function are factored out in the abstract Thin\+Var\+Dep\+Interface class, and some other aspects are factored out in the different abstract Thin\+Compute\+Interface class.

Then, to get to actual concrete objects one only have to define actual Function; currently what is provided is Linear\+Function, which then allows to represent the all-\/important class of Mixed-\/\+Integer Linear Problems. Other concrete functions can be easily defined. Also, the (abstract) class One\+Var\+Constraint is separately defined for Row\+Constraint depending on only one variable (and which therefore don\textquotesingle{}t really need a Function), from of which Box\+Constraint and some specialized versions (like N\+N\+Constraint and N\+P\+Constraint for non-\/negativity and non-\/positivity Constraint, respectively) are derived.

The system is completed by Solver, an abstract base class setting the general interface between a Block representing a mathematical model and any algorithm designed to solve (i.\+e., produce solutions for) it, together with a slightly specialized derived C\+D\+A\+Solver aimed at solution algorithms that exploit convex structures (and, in particular, duality) in the model. Solver also share the Thin\+Compute\+Interface interface, which defines how a (potentially) computationally demanding task is to be executed, and how (potentially, many) parameters controlling it can be provided. The latter is done via a Compute\+Configuration object, deriving from the general Configuration class intended to represent objects for collecting many (structured) algorithmic parameters. The class is also the basis for Block\+Configuration (collecting parameters for different tasks that a Block has to accomplish), and Block\+Solver\+Configuration (collecting parameters for all the Solver attached to a Block and its sub-\/\+Block, recursively). Configuration also offers a factory (to which derived classes have to explicitly register, say using the provided macros).

The idea is that Solver will be specialized for specifically-\/structured mathematical models, corresponding to specific derived classes of Block using specific forms of Variable, Constraint and Objective. Block is meant to represent the general concept of \char`\"{}a part of a mathematical model with a
well-\/understood identity\char`\"{}; that is, derived classes will each represent a model with a specific structure (say, a Knapsack, a Traveling Salesman Problem, a Semi\+Definite program, ...). Hence, specialized Solver will be able to exploit this {\itshape semantically defined} structure (the \char`\"{}physical
representation\char`\"{} of the Block) to provide faster and more effective solution methods, possibly without ever using the \char`\"{}abstract representation\char`\"{} of the Block in terms of its Variable, Constraint and Objective. Solver can be attached to individual sub-\/\+Blocks of a larger, composed Block (recursively), easing the implementation of solution algorithms based on advanced techniques (like decomposition ones). However, each specialized Block will always be able to provide an \char`\"{}abstract representation\char`\"{} of the Block in terms of its Variable, Constraint and Objective, in order to be able to implement solution methods that mix general-\/purpose Solver for large classes of mathematical models (say, Mixed-\/\+Integer Linear Programs) and ad-\/hoc Solver for very specific structures. Block\+Solver\+Configuration is specifically constructed to support the notion that \char`\"{}a Solver for a Block\char`\"{} can actually be a large collection of Solver attached to the Block and its sub-\/\+Block (recursively). In order to allow the Solver to be automatically created and attached to their Block via a Block\+Solver\+Configuration, Solver offers a factory (to which derived classes have to explicitly register, say using the provided macros).

S\+M\+S++ currently supports the following general mechanisms\+:


\begin{DoxyItemize}
\item changes in the data of a Block, Variable, Constraint, Objective, Function ... are treated in a \char`\"{}fully lazy\char`\"{} way, i.\+e., are recorded into objects of (classed derived from the base) class Modification, that are then shipped to all Solver registered to the Block (and all its ancestors).
\item The base Block class supports the \char`\"{}abstract representation\char`\"{} being made of any derived class from Variable, Constraint and Objective, using boost\+::any to support individual objects, arrays and multi-\/dimensional arrays (using boost\+:multi\+\_\+array) of objects.
\item Block provides a factory object to dynamically generate object of any class derived by Block just passing it the string of the classname (provided the class is registered to the factory, for which support is provided with some macros).
\item Block supports (although only in the interface, not in any practical implementation) the concept that a Block can produce different versions of itself that are either equivalent (a Reformulation), a Relaxation or a Restriction. The interface caters for the fact that both solution information (see Solution) and Modification can be mapped back and forth from a Block to any one of its \char`\"{}\+R3 Block\char`\"{}, although this is all predicated on the fact that the Block itself implements this.
\item Variable and Constraint of a Block can be both static and dynamic, in order to allow strategies where they can be dynamically generated while providing the Solver with information about which are never going to change.
\item The solution status of a Block (the value of its Variable) can be saved in appropriate objects derived from the base Solution class, and then read back from these in the Blocl. The Solution objects can possibly save any subset of the information (as dictated by a Configuration object) and dual solution if available. A very general Col\+Variable\+Solution object is provided for doing this only using the \char`\"{}abstract representation\char`\"{} for Block that only have Col\+Variable in them.
\item there is some support for saving an entire Block as a model-\/description file and for loading it from a model-\/description file; currently, methods using the net\+C\+DF library are being tested.
\end{DoxyItemize}

S\+M\+S++ currently {\itshape does not} support a number of important aspects, such as asynchronous execution of the computationally heavy parts (see Thin\+Compute\+Interface). However, it is hopefully structured in such a way that these aspects can be added later with relatively minimal disruption of the current interface.

\section*{Legal Stuff }

\subsection*{Standard Disclaimer }

The code is currently provided free of charge for academic purposes only. As such, it is provided \char`\"{}as is\char`\"{}, without any explicit or implicit warranty that it will properly behave or it will suit your needs. The Authors of the code cannot be considered liable, either directly or indirectly, for any damage or loss that anybody could suffer for having used it. More details about the non-\/warranty attached to this code are available in the license description file.

\subsection*{License }

This code is provided free of charge for academic purposes under the \char`\"{}academic license\char`\"{}\+: see the file doc/academicl.\+txt for further details. Because the code is currently in the early stages of its development we prefer a stricter licensing regime w.\+r.\+t. L\+G\+PL in order to discourage too early fragmentation of the code. Thus, you can make changes, but we strongly suggest that you do not distribute any modified version (although you are legally allowed to within the terms of the license) without prior agreement with the original developers; if your changes make good sense, please allow us to incorporate them in the standard release. Yet, it is foreseen that the license will be moved to L\+G\+PL as soon as the code is deemed stable enough to be widely distributed.

\section*{Version }


\begin{DoxyItemize}
\item Current version of S\+M\+S++ is\+: 0.\+10
\item Current date is\+: 14 -\/ Nov -\/ 2018
\end{DoxyItemize}

\section*{Authors }

\subsection*{Current Lead Authors }

\begin{DoxyVerb}Antonio Frangioni
Operations Research Group
Dipartimento di Informatica
Universita' di Pisa

Rafael Durbano Lobato
Department of Applied Mathematics
State University of Campinas, Brazil
\end{DoxyVerb}


\subsection*{Previous Lead Authors and Contributors }

\begin{DoxyVerb}Kostas Tavlaridis-Gyparakis
Operations Research Group
Dipartimento di Informatica
Universita' di Pisa

Utz-Uwe Haus
Cray EMEA Research Lab\end{DoxyVerb}
 